\section{Conclusion}
\label{sec:conclusion}

We presented \ours, a multi-agent memory system that keeps complete typed
records, indexes them with a provenance hypergraph under explicit budgets, and
runs AutoDream outside the write path to merge duplicate anchors and refine
skills. It leads the compared systems on four of the five QA benchmarks, two of
them with AutoDream disabled, reaching 91.4\% on LoCoMo and 72.6\% on
LongMemEval-S, and it has the highest value on all four metrics for
cross-domain review summarization on MemoryCD. On LoCoMo, the hypergraph and
AutoDream each add accuracy over flat memory. As procedural memory, its
distilled skills lead the compared memory methods in mean success on ALFWorld
and WebShop, and with GPT-5.2 they reach 49.2\% on MetaClaw-Bench, against
44.0\% reported for the MetaClaw skill library. With AutoDream disabled, \ours
answers each question in less time than HyperMem on all five timed benchmarks,
and on WebShop its skills cost the fewest tokens among the memory methods. One
store serves both uses, and the system runs with closed and open models without
fine-tuning. The results suggest that a memory layer can be
practical, accurate, and efficient at once when complete records remain the
system of record, the graph is only a bounded index, and consolidation changes
how evidence is connected rather than what evidence exists. 

\paragraph{Limitations.}
Our ablations add components in groups. They do not isolate n-ary frames,
cited source rows, occurrence identities, answer-time verification, or the
judge, distiller, and curator of procedural memory, so the explanations in \Cref{sec:experiments}
remain interpretations. The user-memory results come from single runs, partly
scored by an LLM judge, the agent results average three runs with large
spreads, and some baselines are values reported by their authors. 

% Building the
% graph slows the ingestion of long documents, the cost of AutoDream is not yet
% measured, and the public benchmarks we use do not test months of real use.
